\section{Centered-moment order}
\label{sec:centered-order}

\begin{discussion}
  \item Explain why the critical profile from
  \cref{def:critical-profile} is the correct center for the generator.
  \item Distinguish centered-moment order from stochastic order.
  \item Explain that the proof below uses a Metzler generator in the
  centered-monomial basis; it does not use conditional patch
  factorization, although both mechanisms originate in
  \patchpositivity.
\end{discussion}

\begin{definition}[Centered-moment order]
\label{def:centered-order}
For finite signed measures \(\lambda_0,\lambda_1\), write
\[
  \lambda_0\preceq_\star\lambda_1
  \quad\Longleftrightarrow\quad
  \lambda_0(\chi_A^\star)
  \le
  \lambda_1(\chi_A^\star)
  \quad\text{for every }A\Subset\Lambda.
\]
\end{definition}

\begin{theorem}[Preservation of centered-moment order]
\label{thm:centered-order}
Let \((P_t)_{t\ge0}\) satisfy
\cref{ass:standing-locality} and have \patchpositivity. If
\(\nu_0\preceq_\star\nu_1\), then, for every
\(A\Subset\Lambda\) and \(t\ge0\),
\begin{align}
  \nu_0(P_t\chi_A^\star)
  &\le
  \nu_1(P_t\chi_A^\star),
  \label{eq:centered-order}
  \\
  \nu_0(P_t\chi_A)
  &\le
  \nu_1(P_t\chi_A).
  \label{eq:ordinary-moment-order}
\end{align}
In particular,
\begin{equation}
  \mu\in\Mstar
  \quad\Longrightarrow\quad
  \mu P_t\in\Mstar.
  \label{eq:Mstar-preserved}
\end{equation}
\end{theorem}

\begin{proof}
For \(S\subseteq N(i)\), define
\[
  d_i(S)
  =
  c_i^0(S)
  -
  p_i^\star\bigl(c_i^0(S)+c_i^1(S)\bigr),
\]
and, for \(S\ne\vn\),
\[
  b_i(S)=-c_i^0(S)-c_i^1(S).
\]
\Cref{thm:patch-positivity-criterion,prop:critical-density-formula}
give
\begin{equation}
  d_i(S)\ge0,
  \qquad
  b_i(S)\ge0.
  \label{eq:db-positive}
\end{equation}
Writing \(r_i=c_i^0(\vn)+c_i^1(\vn)\), a generator calculation gives
\begin{align}
  \cL\chi_A^\star
  =
  \sum_{i\in A}
  \Bigg[
  -r_i\chi_A^\star
  &+
  \sum_{S\subseteq N(i)}
  d_i(S)\chi_S\chi_{A\setminus\{i\}}^\star
  \notag\\
  &+
  \sum_{\vn\ne S\subseteq N(i)}
  b_i(S)\chi_S\chi_A^\star
  \Bigg].
  \label{eq:centered-generator}
\end{align}
For finite \(B,S\),
\begin{equation}
  \chi_S\chi_B^\star
  =
  \left(
    \prod_{j\in S\cap B}(1-p_j^\star)
  \right)
  \sum_{R\subseteq S}
  \left(
    \prod_{j\in S\setminus R}p_j^\star
  \right)
  \chi_{R\cup(B\setminus S)}^\star.
  \label{eq:centered-product-expansion}
\end{equation}
Every coefficient in \eqref{eq:centered-product-expansion} is
nonnegative. In finite volume, the matrix of \(\cL\) in the
centered-monomial basis is therefore Metzler, so its exponential has
nonnegative entries. Finite-range approximation passes the conclusion
to infinite volume.

Apply this positivity-preserving operator to
\(\nu_1-\nu_0\) to obtain \eqref{eq:centered-order}. Finally,
\[
  \chi_A
  =
  \sum_{B\subseteq A}
  \left(
    \prod_{i\in A\setminus B}p_i^\star
  \right)
  \chi_B^\star
\]
has nonnegative coefficients, which gives
\eqref{eq:ordinary-moment-order}. The defining inequalities of
\(\Mstar\) then give \eqref{eq:Mstar-preserved}.
\end{proof}

\begin{corollary}[Product profiles]
\label{cor:product-profile-order}
If
\[
  \pstar\le\mathbf p\le\mathbf q,
\]
then
\[
  \mu_{\mathbf p}(P_t\chi_A)
  \le
  \mu_{\mathbf q}(P_t\chi_A).
\]
The same comparison holds for mixtures whose random profiles admit
such a coordinatewise coupling.
\end{corollary}

\begin{proof}
For every finite \(B\),
\[
  \mu_{\mathbf p}(\chi_B^\star)
  =
  \prod_{i\in B}(p_i-p_i^\star)
  \le
  \prod_{i\in B}(q_i-p_i^\star)
  =
  \mu_{\mathbf q}(\chi_B^\star).
\]
Apply \cref{thm:centered-order}, then average under the coupling.
\end{proof}

\begin{corollary}[Reflected product comparison]
\label{cor:reflected-product-comparison}
If
\[
  \mathbf q\le\mathbf p,
  \qquad
  \mathbf q+\mathbf p\ge2\pstar,
\]
then
\[
  \mu_{\mathbf q}(P_t\chi_A)
  \le
  \mu_{\mathbf p}(P_t\chi_A).
\]
\end{corollary}

\begin{proof}
The assumptions give
\(\abs{q_i-p_i^\star}\le p_i-p_i^\star\), so every centered moment of
\(\mu_{\mathbf q}\) is bounded above by the corresponding centered
moment of \(\mu_{\mathbf p}\). Apply
\cref{thm:centered-order}.
\end{proof}

\begin{proposition}[Invariant measure in a preserved class]
\label{prop:invariant-preserved-class}
Let \(P_t\) be Feller on a compact configuration space. Every nonempty,
weakly closed, convex class of probability measures preserved by
\((P_t)\) contains an \invariantmeasure.
\end{proposition}

\begin{proof}
For a measure \(\mu\) in the class, form the Cesàro averages
\[
  \mu_T=\frac1T\int_0^T\mu P_s\,\dd s.
\]
They remain in the class. Compactness gives a weakly convergent
subsequence. The semigroup property and Feller property show that its
limit is invariant. This is the Krylov--Bogoliubov argument; see
\citet{liggett1985}.
\end{proof}

\begin{corollary}
\label{cor:invariant-in-Mstar}
Under the Feller assumption, \(\Mstar\) contains an
\invariantmeasure.
\end{corollary}

\begin{proof}
\Cref{prop:centered-classes} makes \(\Mstar\) nonempty, convex, and
weakly closed, while \cref{thm:centered-order} shows it is preserved.
Apply \cref{prop:invariant-preserved-class}.
\end{proof}
